\SourcePage{231}{236}
\section{Radiation from diatomic molecules. Vibrational and rotational spectra}

The selection rules and transition-probability formulas listed in the preceding section remain valid for transitions in which the molecule's electronic state does not change\footnote{Transitions that alter the vibrational state (and, with it, the rotational state) constitute what is called the molecule's \emph{vibrational spectrum}; it lies in the near-infrared region (wavelengths below 20~$\mu$m). Transitions that alter only the rotational state constitute the \emph{rotational spectrum}, which lies in the far-infrared region (wavelengths above 20~$\mu$m).}. We shall consider here only a few features particular to these transitions.

First of all, selection rule (53.4) forbids dipole transitions without a change of electronic state in molecules made of identical atoms: such a transition would leave the parity of the electronic term unchanged. As follows from \S~53, this prohibition can be violated only when the interaction of the nuclear spins with the electrons is taken into account; in molecules composed of different isotopes of one element, the effect of rotation on the electronic state is already sufficient to violate it.

The calculation of dipole-moment matrix elements reduces (by the formulas of \S~87 (see III)) to calculating them in a coordinate system that rotates with the molecule. In this system the molecular wave function is the product of the electron wave function at a specified separation $r$ of the nuclei and the wave function for the vibrational motion of the nuclei in the effective field $U(r)$ of the electrons and nuclei. If the influence of nuclear motion on the electronic state is neglected completely, the initial and final electron wave functions for the transitions under consideration are identical. Integration over the electron coordinates therefore leaves in the matrix element simply the mean dipole moment $\bar{\mathbf d}$ of the molecule (directed, evidently, along its axis), as a function of the separation $r$. Since the vibrations are small, $\bar{\mathbf d}(r)$ can be expanded in powers of the vibrational coordinate $q=r-r_0$. For transitions involving a change of vibrational state, the constant term of this expansion drops out of the matrix element because the wave functions for vibrational motion in the same field $U(q)$ are orthogonal; the term proportional to $q$ remains. If the vibrations are treated as harmonic, the familiar properties of the
\SourcePage{232}{237}
linear oscillator (see III, \S~23) show that the matrix elements are nonzero only for transitions between adjacent vibrational states. Thus the vibrational quantum number $v$ obeys the selection rule
\begin{equation}
v'-v=\pm1.
\tag{54.1}
\end{equation}
This rule is violated, however, when vibrational anharmonicity and the subsequent terms in the expansion of $\bar{\mathbf d}(q)$ are included.

For a purely rotational transition (with no change of vibrational state either), the matrix element of the dipole-moment projection on the moving axis $\zeta$ may be replaced simply by the mean molecular dipole moment $\bar{\mathbf d}=\bar{\mathbf d}(0)$\footnote{For a molecule made of identical atoms, $\bar{\mathbf d}=0$, as is evident from symmetry.}. The resulting formula for the probability of the transition $J\to J-1$ is
\begin{equation}
w(nJ\to n,J-1)=\frac{4\omega^3}{3\hbar c^3}\bar{\mathbf d}^{\,2}\frac{J^2-\Omega^2}{J(2J+1)},
\tag{54.2}
\end{equation}
which gives not only relative probabilities, as does (53.12), but their absolute values as well. (Formula (54.2) is written for case $a$; in case $b$, $K,\Lambda$ must replace $J,\Omega$.)

The frequencies of purely rotational transitions are the differences between the rotational energies $BJ(J+1)$ and are therefore
\begin{equation}
\hbar\omega_{J,J-1}=2BJ.
\tag{54.3}
\end{equation}
Successive lines are separated by the same interval, $2B$.
